\subsubsection{Monitoring instrument performance}

The Universit\'e de Montr\'eal group has developed and sustained a strong expertise in monitoring \spi\ data quality and science readiness since the instrument's first light. This work has included systematic checks of the RV zero-point, detector drifts, wavelength-solution stability, flux sensitivity, and the overall integrity of calibration sequences, and it has directly informed several calibration-strategy updates implemented in \texttt{APERO} \citep{cook_apero_2022}.

For \amakihi, we will extend and formalize this effort through the \texttt{APERO} Reduction Interface (ARI), a web-based dashboard that provides rapid access to quality metrics, observation summaries, and contextual diagnostics for each visit. ARI will allow the team and CFHT staff to track both engineering-level and science-level indicators in a homogeneous framework, helping to identify trends, diagnose issues early, and document the long-term behavior of the instrument.

At the engineering level, ARI will monitor the main quality-control quantities delivered by the reduction chain, including image-shape diagnostics, drift measurements, cavity behavior, and wavelength-solution centroids. At the science level, it will track quantities such as signal-to-noise ratio, barycentric velocity coverage, telluric indicators, saturation flags, readout noise, reduction version, processing date, calibration age, and the status of LBL and CCF products. The interface will also provide direct access to reduced data products and to reduction logs, allowing users to identify what was processed, what failed, and under which calibration context.

This monitoring framework is intended both as an internal operational tool and as a community-facing resource. In practice, it will support night-to-night assessment of instrument health, rapid validation of newly acquired observations, and long-term cross-checks of the stability required for precision-RV science. By consolidating these diagnostics in a single interface, ARI will strengthen the reproducibility of the reductions, improve responsiveness to instrumental or environmental anomalies, and help ensure that the data delivered by \wena\ remain uniformly science-ready throughout the survey.
